\documentclass[conference]{IEEEtran}
\IEEEoverridecommandlockouts

\usepackage{cite}
\usepackage{amsmath,amssymb,amsfonts}
\usepackage{algorithmic}
\usepackage{graphicx}
\usepackage{textcomp}
\usepackage{xcolor}
\usepackage{listings} % Added FOR Haskell FORmatting
\usepackage{float}
\usepackage{algorithm}
\usepackage{booktabs}
\usepackage{multirow}

\def\BibTeX{{\rm B\kern-.05em{\sc i\kern-.025em b}\kern-.08em
    T\kern-.1667em\lower.7ex\hbox{E}\kern-.125emX}}

\begin{document}

\title{Particle Swarm Optimisation for Large-Scale Optimisation}

\author{\IEEEauthorblockN{Jordan Michael Ward}
\IEEEauthorblockA{\textit{Computer Science Department} \\
\textit{Stellenbosch University}\\
Stellenbosch, South Africa \\
% Email: <insert email here>
}
}

\maketitle
\thispagestyle{plain}
\pagestyle{plain}

\begin{abstract}
    Large scale optimisation problems (LSOPs) present a fundamental challenge for traditional optimisation algorithms due to 
    the curse of dimensionality, a phenomenon exponentially expanding the search space and leading to 
    velocity explosion and premature stagnation. In this report a hybrid particle swarm optimisation (HPSO), 
    integrating both Subspace Initialisation and Stochastic Scaling (two techniques used to better scale LSOPs) is introduced and measured under the hypothesis that the algorithm 
    will outperform standard algorithms on complex multi-dimensional landscapes. By combining both Stochastic Scaling and 
    Subspace Initialisation, it is theorized an increase in performance will be attained when comparing the results of the algorithm 
    to implementations on LSOPs using both approaches separably. Experimental evaluation was performed with dimensionalities 
    100, 500 and 1000 per configuration. Every algorithm ran for 30 independent runs at 5000 iterations per run.
\end{abstract}

\section{Introduction}
    Large Scale Optimisation remains a difficult area in optimisation due to the "Curse of Dimensionality". When the 
    dimension for a problem increases, the search landscape grows exponentially. To help mitigate this scalability issue, approaches such as 
    Subspace Initialisation and Stochastic Scaling have been developed to better navigate high-dimensional problems. The 
    goal of this report is to explore an advanced hybrid particle swarm optimisation (HPSO) model that merges Subspace 
    Initialisation and Stochastic Scaling, and to measure the performance of the model against the algorithms the hybrid model is built from. 

    Four PSO variants were implemented: Canonical, Subspace Initialisation, Stochastic Scaling, and Hybrid. Each 
    algorithm was rigorously benchmarked against five classic mathematical landscapes: Sphere, Elliptic, Rosenbrock, 
    Rastrigin, and Schwefel 1.2. For every benchmark function, 30 independent runs of 5000 iterations were performed 
    at dimensions 100, 500, and 1000 to ensure statistical significance. 

    The empirical results demonstrate that the proposed Hybrid PSO significantly outperforms both the Canonical 
    baseline and the independent constituent algorithms of the baseline on the majority of the tested landscapes. By effectively 
    mitigating high-dimensional velocity explosions while maintaining search discipline, the HPSO achieves 
    stable convergence. However, a stagnation observed on the Rosenbrock function 
    highlights an algorithmic trade-off, revealing that the very mechanisms used to address steep landscapes 
    can become detrimental on flat-valley terrains.

    The remainder of this report is structured as follows. Section 2 introduces the background theory and foundational 
    concepts of PSO. Section 3 details the methodology and mathematical formulations of the advanced variants. Section 
    4 outlines the empirical procedure and benchmark configurations. Section 5 presents the analysis of the 
    results. Finally, Section 6 provides concluding remarks.

\section{Background}
    Global optimisation focuses on finding the best solution, the global maximum or minimum, for a given objective function over a 
    continuous search space. Unlike local search algorithms, where the search stops in the nearest valley or peak, global optimisation 
    aims to navigate complex landscapes to find the single best solution among many potential local optima. A large scale optimisation problem 
    (LSOP) is a global optimisation problem with a high dimensionality, typically in the hundreds or thousands. This phenomenon is the infamous 
    "Curse of Dimensionality" in LSOPs. When the dimension in optimisation parameters increase, the search space grows exponentially,  
    making it impossible to scour the entire space for the global optimum. The exponential growth makes exhaustive searches infeasible, and causes distance 
    metrics to be drastically skewed. 
    
    Because deterministic or exhaustive approaches fail under these extreme conditions, researchers rely on 
    meta-heuristics, like population based evolutionary algorithms, to find "good enough" solutions in a reasonable timeframe. A 
    prominent approach is particle swarm optimisation (PSO), originally introduced by Kennedy and Eberhart [CITATION]. Inspired by the social 
    foraging behaviour of bird flocks, PSO maintains a population, or "swarm", of candidate solutions called particles.
    As these particles move through the multi-dimensional search space, the trajectories of the particles are guided by three main forces: the current momentum of the particles, 
    the historical personal best position, and the global best position found by the entire swarm. 

    The algorithm updates the velocity and position of each particle $i$ at iteration $t$ as
    
    \begin{equation}
        V_{i}^{t+1} = w V_{i}^t + c_1 r_1 (Pbest_{i}^t - X_{i}^t) + c_2 r_2 (Gbest^t - X_{i}^t)
        \label{eq:pso_velocity}
    \end{equation}
    and
    \begin{equation}
        X_{i}^{t+1} = X_{i}^t + V_{i}^{t+1}
        \label{eq:pso_position}
    \end{equation}
    respectively. In equation \ref{eq:pso_velocity}, $w$ represents the inertia weight of the particle [CITATION]; the inertia weight controls the momentum of the particle. 
    The cognitive and social acceleration coefficients are represented by $c_1$ and $c_2$; the coefficients control how strongly the particle is pulled 
    back to its own personal historic best ($Pbest$), or the global best of the swarm ($Gbest$).
    Finally, $r_1$ and $r_2$ are stochastic vectors sampled uniformly from $[0, 1]^D$, introducing the stochastic scaling used to 
    encourage diversity and explore the search space. 

    The general framework of the canonical PSO is outlined in Algorithm \ref{alg:canonical_pso}.

    \begin{algorithm}[H]
    \caption{Canonical Particle Swarm Optimisation}
    \label{alg:canonical_pso}
    \begin{algorithmic}[1] % [1] adds line numbers
    \STATE \textbf{Initialise} swarm with $N$ particles
    \FOR{each particle $i \in 1..N$}
        \STATE Randomly initialise position $X_i$ within search bounds
        \STATE Initialise velocity $V_i$ to randomly scaled small magnitudes
        \STATE Evaluate objective fitness $f(X_i)$
        \STATE Set personal best $Pbest_i = X_i$
    \ENDFOR
    \STATE Determine global best $Gbest$ from all $Pbest_i$
    \WHILE{termination condition not met}
        \FOR{each particle $i \in 1..N$}
            \STATE Update velocity $V_i$ using Equation \ref{eq:pso_velocity}
            \STATE Update position $X_i$ using Equation \ref{eq:pso_position}
            \STATE enforce boundary constraints on $X_i$
            \STATE Evaluate new objective fitness $f(X_i)$
            \IF{$f(X_i)$ is better than $f(Pbest_i)$}
                \STATE Update $Pbest_i = X_i$
            \ENDIF
        \ENDFOR
        \STATE Update $Gbest$ if any new $Pbest_i$ is better than current $Gbest$
    \ENDWHILE
    \RETURN $Gbest$
    \end{algorithmic}
    \end{algorithm}

    While canonical PSO performs well on low-dimensional spaces, its velocity updates and uniform random initialisation often lead to premature 
    convergence or stagnation when subjected to the sheer volume of Large Scale Optimisation Problems [CITATION]. To address this challenge, modern 
    literature proposes modifying the standard initialisation and velocity mechanisms. 
\section{Methodology}
    This section details the algorithm architectures to test against the chosen 
    Large Scale Optimisation Problems. The algorithms were implemented in Python, using the cilpy 
    library. 

\subsection{Standard Particle Swarm Optimisation}
    The Canonical PSO algorithm acts as the baseline in this experiment. It follows the standardised formulation 
    and uniform random initialisation pattern defined in Algorithm \ref{alg:canonical_pso}, ensuring a standard benchmark 
    providing a baseline to measure more advanced scaling methods.

\subsection{Stochastic Scaling Particle Swarm Optimisation}
    In the Canonical PSO, unique random vectors $r_1$ and $r_2$ are applied to every single dimension to encourage diversity 
    and inject random noise. In high-dimensional problems, however, this independent scaling model actually causes the 
    swarm to behave erratically, with a high degree of roaming. To mitigate this, the dimensions for the swarm are randomly partitioned into $k$ 
    sub groups, a technique known as Stochastic Scaling. 

    During every iteration, all dimensions are 
    randomly shuffled into $k$ equal-sized groups, where a single $r_1$ and $r_2$ value is then drawn and multiplied 
    with every dimension within its respective group. The value of $k$ is not static; it starts at one for iteration zero and increases linearly 
    up to the maximum dimensionality $D$.  

\subsection{Subspace Initialisation Particle Swarm Optimisation}
    Particles are uniformly initialised over the entire search space in the Canonical PSO, in an attempt to acquire a broad understanding of the entire
    landscape. Unfortunately, in high-dimensional spaces, this initialisation model leads to chaotic velocity explosion, causing particles to overshoot desirable regions and bounce 
    off boundary walls. Subspace Initialisation attempts to counter this effect by modifying how the swarm particles are initially distributed, severely 
    constricting the swarms starting positions to a heavily constrained area. The algorithm calculates the geometric center of the search space, and 
    calculates a one-dimensional line that passes directly through the origin. The swarm is then placed all along this line with the initial velocities of the particles 
    set to zero. By initialising the swarm in this manner the algorithm promotes fine-grained exploitation and mitigates 
    velocity explosion, before the particles begin to fan out and explore wider spaces. 

\subsection{Hybrid Particle Swarm Optimisation}
    The Hybrid PSO, as seen in Algorithm \ref{alg:hybrid_pso}, combines both aforementioned scaling techniques for LSOPs into a single algorithmic structure. The goal of the hybrid model is to harness the 
    benefits of the structured starting positions of Subspace Initialisation, which mitigates velocity explosion at the start, with the dynamic and randomised 
    sub-group velocity updates of Stochastic Scaling. This combination controls particle roaming as well as massive step sizes throughout the experiment. 
    
    The swarm is initially seeded entirely using the Subspace approach, providing a structural exploration footing. Then, during the iterative 
    loops, the velocity equations are driven entirely by the linearly scaling $k$-group dimensionality of the Stochastic model. It was hypothesised that 
    the hybrid structure will outperform both independent algorithms tested individually. 

    \begin{algorithm}
    \caption{Hybrid Particle Swarm Optimisation (Subspace Initialisation + Stochastic Scaling)}
    \label{alg:hybrid_pso}
    \begin{algorithmic}[1]
    \STATE \textbf{Phase 1: Subspace Initialisation}
    \STATE Calculate the geometric center $C$ of the search space bounds
    \STATE Generate an orthonormal seed vector $d$ via Modified Gram-Schmidt
    \FOR{each particle $i \in \{1, \dots, N\}$}
        \STATE Generate position $p_1$ along the line $p = t \cdot d + C$ (with margin)
        \STATE Generate position $p_2$ along the line $p = t \cdot d + C$ (strict bounds)
        \IF{$f(p_2) < f(p_1)$}
            \STATE Personal best $y_i \gets p_2$, Initial position $x_i \gets p_1$
        \ELSE
            \STATE Personal best $y_i \gets p_1$, Initial position $x_i \gets p_2$
        \ENDIF
        \STATE Initial velocity $v_i \gets \mathbf{0}$ \COMMENT{Eliminate initial momentum}
    \ENDFOR
    \STATE Initialise global best $\hat{y}$
    \STATE
    \STATE \textbf{Phase 2: Stochastic Scaling Loop}
    \FOR{iteration $t = 1$ to $T_{max}$}
        \STATE Calculate group count $k$ (linearly increasing from $1$ to $D$)
        \STATE Randomly partition the $D$ dimensions into $k$ groups
        \FOR{each group $g \in \{1, \dots, k\}$}
            \STATE Draw shared group scalars $r_{1g} \sim U(0,1)$ and $r_{2g} \sim U(0,1)$
        \ENDFOR
        \FOR{each particle $i \in \{1, \dots, N\}$}
            \FOR{each dimension $d \in \{1, \dots, D\}$}
                \STATE Let $g$ be the specific group containing dimension $d$
                \STATE $v_{id} \gets w v_{id} + c_1 r_{1g} (y_{id} - x_{id}) + c_2 r_{2g} (\hat{y}_d - x_{id})$
                \STATE $x_{id} \gets x_{id} + v_{id}$
                \STATE Clamp $x_{id}$ to search space bounds
            \ENDFOR
            \STATE Evaluate $f(x_i)$ and update personal best $y_i$
        \ENDFOR
        \STATE Update global best $\hat{y}$ if necessary
    \ENDFOR
    \end{algorithmic}
    \end{algorithm}

\section{Empirical Procedure}
    To evaluate the algorithms, the experimental setup was designed to stress test the Particle Swarm Optimisation 
    variants against increasingly large search spaces to capture how algorithmic performance degenerates or improves with dimensionality. 

\subsection{Benchmark Functions}
    The capability of the algorithm variants was tested against five classic benchmarking optimisation functions, consisting of both 
    unimodal and multi-modal landscapes. The functions were evaluated across dimensions ($D$) $100$, $500$, and $1000$, and the specific character domains 
    were as follows:
    
    \begin{itemize}
        \item \textbf{Sphere}: A smooth, unimodal, strongly convex landscape. Search domain: $[-100.0, 100.0]^D$.
        \item \textbf{Elliptic}: A highly ill-conditioned function, representing stretched topologies. Search domain: $[-100.0, 100.0]^D$.
        \item \textbf{Rastrigin}: A highly multi-modal function containing numerous periodic local minima. Search domain: $[-5.12, 5.12]^D$.
        \item \textbf{Rosenbrock}: A non-convex, narrow "valley" topology testing local minima avoidance. Search domain: $[-30.0, 30.0]^D$.
        \item \textbf{Schwefel 1.2}: A deeply non-separable unimodal function. Search domain: $[-100.0, 100.0]^D$.
    \end{itemize}

\subsection{Experimental Parameters}
    Fairness across all algorithms was strictly maintained. Every algorithm executed for exactly $30$ independent runs to ensure robust 
    averages. Each single run was strictly terminated after $5000$ iterations. 
    The swarm size was tightly constrained to $30$ particles, the inertia weight ($w$) was set to $0.729844$, and both the 
    cognitive ($c_1$) and social ($c_2$) acceleration coefficients were defined equally as $1.49618$.
    Because standard baseline parameters across every algorithm were used, any performance discrepancies measured are attributed to the 
    underlying algorithm architecture.

\subsection{Evaluation Metrics}
    To determine the efficiency and accuracy of the underlying architectures, final fitness values across all benchmarking functions were logged 
    for all $30$ independent runs. The \textbf{Mean Error} and \textbf{Standard Deviation} are reported under tabular conditions in Table \ref{tab:results}.
    
    Additionally, logarithmic convergence profiles plotting the generation sequence against the global best fitness evaluations are provided to 
    visually represent optimisation speed and trajectory stagnation characteristics over time.

\begin{table*}[htbp]
\centering
\caption{Mean Best Fitness and Standard Deviation over 30 Independent Runs}
\label{tab:results}
\begin{tabular}{llcccc}
\toprule
\textbf{Function} & \textbf{$D$} & \textbf{Canonical PSO} & \textbf{Stochastic Scaling} & \textbf{Subspace Init} & \textbf{Hybrid PSO} \\
\midrule
\multirow{3}{*}{\textbf{Sphere}} & 100 & 1.43e+04 $\pm$ 8.98e+03 & 3.33e+02 $\pm$ 1.83e+03 & 7.55e-08 $\pm$ 2.93e-07 & \textbf{2.73e-11 $\pm$ 1.44e-10} \\
 & 500 & 2.82e+05 $\pm$ 3.86e+04 & 2.23e+04 $\pm$ 7.48e+03 & 1.81e+01 $\pm$ 3.19e+01 & \textbf{4.40e-01 $\pm$ 7.48e-01} \\
 & 1000 & 1.02e+06 $\pm$ 6.98e+04 & 1.56e+05 $\pm$ 3.11e+04 & 5.34e+01 $\pm$ 9.31e+01 & \textbf{4.98e+00 $\pm$ 6.43e+00} \\
\midrule
\multirow{3}{*}{\textbf{Elliptic}} & 100 & 8.59e+08 $\pm$ 5.15e+08 & 1.99e+07 $\pm$ 6.11e+07 & 2.54e-04 $\pm$ 9.82e-04 & \textbf{1.38e-09 $\pm$ 4.91e-09} \\
 & 500 & 1.31e+10 $\pm$ 3.99e+09 & 3.78e+08 $\pm$ 1.89e+08 & 2.89e+05 $\pm$ 8.77e+05 & \textbf{4.94e+03 $\pm$ 8.96e+03} \\
 & 1000 & 3.29e+10 $\pm$ 6.32e+09 & 2.89e+09 $\pm$ 5.75e+08 & 1.58e+06 $\pm$ 2.37e+06 & \textbf{2.52e+05 $\pm$ 4.14e+05} \\
\midrule
\multirow{3}{*}{\textbf{Rosenbrock}} & 100 & 2.40e+07 $\pm$ 4.28e+07 & 2.01e+02 $\pm$ 5.67e+01 & 8.89e+01 $\pm$ 1.93e+00 & \textbf{8.73e+01 $\pm$ 9.90e-01} \\
 & 500 & 1.23e+09 $\pm$ 2.69e+08 & 7.65e+06 $\pm$ 3.06e+06 & 6.04e+02 $\pm$ 2.87e+02 & \textbf{4.99e+02 $\pm$ 1.16e+01} \\
 & 1000 & 4.44e+09 $\pm$ 4.63e+08 & 8.52e+07 $\pm$ 3.33e+07 & 1.92e+03 $\pm$ 2.06e+03 & \textbf{1.03e+03 $\pm$ 8.94e+01} \\
\midrule
\multirow{3}{*}{\textbf{Rastrigin}} & 100 & 6.21e+02 $\pm$ 8.48e+01 & 1.88e+02 $\pm$ 3.67e+01 & 6.30e-01 $\pm$ 2.42e+00 & \textbf{8.69e-09 $\pm$ 2.45e-08} \\
 & 500 & 3.96e+03 $\pm$ 1.64e+02 & 1.58e+03 $\pm$ 1.46e+02 & 1.02e+01 $\pm$ 2.36e+01 & \textbf{3.06e-01 $\pm$ 4.75e-01} \\
 & 1000 & 9.49e+03 $\pm$ 2.30e+02 & 5.10e+03 $\pm$ 2.66e+02 & 2.77e+01 $\pm$ 3.79e+01 & \textbf{6.20e+00 $\pm$ 1.97e+01} \\
\midrule
\multirow{3}{*}{\textbf{Schwefel12}} & 100 & 8.93e+04 $\pm$ 4.32e+04 & 1.02e+04 $\pm$ 1.38e+04 & 3.94e+00 $\pm$ 6.35e+00 & \textbf{1.65e-03 $\pm$ 7.03e-03} \\
 & 500 & 1.77e+06 $\pm$ 3.79e+05 & 7.22e+05 $\pm$ 2.69e+05 & 1.77e+03 $\pm$ 3.86e+03 & \textbf{7.61e+01 $\pm$ 1.70e+02} \\
 & 1000 & 6.89e+06 $\pm$ 9.79e+05 & 3.07e+06 $\pm$ 9.26e+05 & 7.73e+03 $\pm$ 2.47e+04 & \textbf{2.67e+02 $\pm$ 4.08e+02} \\
\bottomrule
\end{tabular}
\end{table*}


\section{Research Results}
This section presents the empirical findings of the study. First, convergence behaviours 
and structural biases are analysed. Following this, a specific note on the Rosenbrock anomaly is discussed.

\subsection{Convergence Behaviours and Algorithmic Structural Bias}
The convergence profiles for $D=1000$ across the benchmark functions are presented in Figures \ref{fig:sphere1000} 
through \ref{fig:rosenbrock1000}. Across the Sphere, Elliptic, Rastrigin, and Schwefel 1.2 landscapes, a consistent 
performance hierarchy emerges.

The Canonical PSO consistently exhibited the poorest performance. Driven by independent, per-dimension random 
initialisation, the Canonical swarm suffers from severe velocity explosion in high-dimensional spaces, causing 
particles to rapidly expand the search radius of the particles and waste evaluations on out-of-bounds roaming. 

The purely Stochastic Scaling PSO successfully dampens this velocity explosion through dimensional 
coupling, keeping the swarm stable. However, the stochastic scaling algorithm plateaus early in the search process. Because the initial positions of the stochastic scaling swarm 
are scattered uniformly at random, the coupled velocity updates restrict the reachable space of the swarm, 
preventing the swarm from navigating to the global optimum from a distant starting point. 

The Subspace Initialisation PSO bypasses the initial velocity explosion by seeding the swarm entirely along a 
one-dimensional line passing through the center of the bounds. It is important to note the global optima of these standard benchmark 
functions lie exactly at the origin. Therefore, Subspace 
Initialisation effectively drops the swarm precisely on the target landscape at Iteration 0, providing a massive, 
head-start. However, because the algorithm subsequently reverts to standard Canonical PSO velocity updates, 
the subspace initialisation algorithm eventually succumbs to high-dimensional chaos and struggles to fine-tune the exploitation of the swarm.

The \textbf{Hybrid PSO} heavily dominates these four benchmarks, but the superiority of the hybrid model is not solely due to this 
initialisation advantage. While the hybrid PSO shares the exact same origin-drop advantage as the Subspace algorithm, the Hybrid 
immediately transitions into Stochastic Scaling for its iterative velocity updates. This mechanism actively controls 
particle momentum, preventing the swarm from exploding off the initial subspace line into chaotic jitter. This 
biased starting position with damped velocity mechanics results in a 
staircase convergence, allowing the swarm to converge 
significantly deeper into the local minima than pure Subspace Initialisation.

\begin{figure}[htbp]
    \centering
    \includegraphics[width=\linewidth]{../plots/Sphere_D1000.png}
    \caption{Convergence Profile for Sphere ($D = 1000$)}
    \label{fig:sphere1000}
\end{figure}

\begin{figure}[htbp]
    \centering
    \includegraphics[width=\linewidth]{../plots/Elliptic_D1000.png}
    \caption{Convergence Profile for Elliptic ($D = 1000$)}
    \label{fig:elliptic1000}
\end{figure}

\begin{figure}[htbp]
    \centering
    \includegraphics[width=\linewidth]{../plots/Schwefel12_D1000.png}
    \caption{Convergence Profile for Schwefel 1.2 ($D = 1000$)}
    \label{fig:schwefel1000}
\end{figure}

\begin{figure}[htbp]
    \centering
    \includegraphics[width=\linewidth]{../plots/Rastrigin_D1000.png}
    \caption{Convergence Profile for Rastrigin ($D = 1000$)}
    \label{fig:rastrigin1000}
\end{figure}

\subsection*{Note on the Rosenbrock Anomaly}

While the Hybrid PSO dominates symmetrically centered topologies, the algorithm immediately stagnates at roughly $10^3$ on the 
Rosenbrock benchmark as illustrated in Figure \ref{fig:rosenbrock1000}. This occurs due to how the algorithm interacts with the notoriously flat parabolic valley of the Rosenbrock function. Subspace Initialisation forces the swarm to start at the origin $[0, \dots, 0]$ with 
zero initial velocity. Trapped on this flat valley floor without initial momentum, and heavily restricted by the 
dampening of Stochastic Scaling, the particles cannot generate the velocity required to navigate toward the global 
minimum at $[1, \dots, 1]$. This highlights a critical trade-off: mechanisms to prevent velocity 
explosions on steep topologies can actively hinder exploration on flat-valley landscapes.

\begin{figure}[htbp]
    \centering
    \includegraphics[width=\linewidth]{../plots/Rosenbrock_D1000.png}
    \caption{Convergence Profile for Rosenbrock ($D = 1000$)}
    \label{fig:rosenbrock1000}
\end{figure}

\section{Conclusion}
The primary objective of this study was to evaluate strategies for mitigating the curse of dimensionality in large scale 
optimisation problems (LSOPs). A hybrid particle swarm optimisation (HPSO) algorithm was developed, integrating 
Subspace Initialisation and Stochastic Scaling. The performance of this hybrid was benchmarked against 
Canonical PSO and the independent constituent algorithms of the hybrid across five landscapes at dimensionalities 
up to $D=1000$.

The empirical results supported the initial hypothesis for the majority of the tested landscapes. By combining 
the geometrically constrained, zero velocity starting conditions of Subspace Initialisation with the  
coupled velocity dampening of Stochastic Scaling, the Hybrid PSO successfully prevented high-dimensional velocity 
explosion and excessive particle roaming. The Hybrid PSO outperformed the Canonical, 
purely Stochastic, and purely Subspace swarms on the Sphere, Elliptic, Rastrigin, and Schwefel 1.2 benchmarks, 
demonstrating stable convergence.

However, this investigation also identified a critical algorithmic limitation through the Rosenbrock anomaly. The 
complete stagnation of the Hybrid swarm on the Rosenbrock function proved that while eliminating initial momentum is 
highly beneficial for symmetrically steep functions, this elimination is actively detrimental on flat-valley landscapes 
where initial momentum is strictly required for traversal. 

Conducting this research highlighted that algorithmic performance in LSOPs cannot be evaluated purely on final 
fitness metrics. An understanding of how the mechanics of an algorithm interact with the specific spatial landscape of a problem is 
also important. 

Future work should address the initialisation bias present in this study by benchmarking the Hybrid 
algorithm against "shifted" landscape variants, where the global optimum is intentionally displaced from the geometric 
center of the search space. Future studies should investigate adaptive initialisation strategies that can  
detect flat-valley landscapes in order to inject necessary initial momentum to address stagnation 
vulnerabilities observed in the Rosenbrock landscape.

\end{document}
